% !TeX encoding = UTF-8
% Build with: python3 build_manual.py
% PDF名の版番号は documentversion から自動取得する。
\documentclass[a4paper,11pt]{ltjsarticle}
\usepackage[margin=23mm]{geometry}
\usepackage[haranoaji]{luatexja-preset}
\usepackage{xcolor}
\usepackage{fvextra}
\usepackage[unicode,hidelinks]{hyperref}
\DefineVerbatimEnvironment{command}{Verbatim}{breaklines=true,breakanywhere=true,fontsize=\small,frame=single,rulecolor=\color{gray},framesep=3mm}
\setlength{\parindent}{0pt}
\setlength{\parskip}{0.45em}
% 更新時は版番号・更新日を変更し、下の更新履歴に行を追加する。
% 内容の修正・追記では小数部を、大きな構成変更では整数部を更新する。
\newcommand{\documentversion}{2.1}
\newcommand{\documentupdated}{2026年10月26日}
\title{MuSEUM実験\\計測・制御プログラム操作マニュアル\\[0.5em]\large 第\documentversion 版}
\author{}
\date{最終更新：\documentupdated}
\begin{document}
\maketitle
\section*{更新履歴}
\begin{tabular}{@{}llp{0.58\linewidth}@{}}
\hline
版番号 & 更新日 & 変更内容 \\
\hline
2.0 & 2026年10月26日 & LaTeX版を作成。 \\
2.1 & 2026年10月26日 & 表題と更新履歴の記載を修正。 \\
\hline
\end{tabular}
\clearpage
\tableofcontents
\clearpage


\section{スイッチの切り替え：Switchcontrol.py}

\subsection{特性評価/実験、TM110/TM210を切り替える}

\begin{command}

% python Switchcontrol.py AMP TM110

\end{command}

\begin{itemize}

\item 引数1 VNA or AMP

\item 引数2 TM110 or TM210

\end{itemize}

\subsection{スイッチ用電源のCH2・CH3・CH4の出力を全てOFFにする}

\begin{command}

% python Switchcontrol.py OFF

\end{command}

\subsection{python起動後}

\begin{command}

>>>import Switchcontrol

>>>Switchcontrol.set_switch("VNA", "TM110")

\end{command}

\begin{itemize}

\item 引数1 VNA or AMP

\item 引数2 TM110 or TM210

\end{itemize}

または


\begin{command}

>>>Switchcontrol.all_off()

\end{command}

\section{VNAの設定・測定・CSV保存：test\_main.py}

\subsection{VNA接続 → 測定条件の設定 → 全4chを各10回測定 → CSV保存をまとめて行う}

\begin{command}

% python test_main.py

\end{command}

\begin{itemize}

\item 保存ファイル名の末尾を指定する

\begin{command}

% python test_main.py TM110_check

\end{command}

\end{itemize}

引数：ファイル名の末尾に付ける文字列


\section{VNAの個別操作：vna\_tools.py}

\subsection{Python起動後、VNAと接続}

\begin{command}

>>> import pyvisa

>>> import vna_tools

>>> rm = pyvisa.ResourceManager('@py')

>>> znb = vna_tools.get_vna_resource(rm, 'TCPIP0::192.168.12.4::inst0::INSTR')

\end{command}

\subsection{測定条件・トレース・画面・マーカーを設定する}

\begin{command}

>>> vna_tools.setup_vna(znb)

\end{command}

\begin{itemize}

\item TM110/TM210それぞれのWide/Narrowを全4chに設定する。

\end{itemize}

\subsection{全4chを各10回測定し、CSVを保存する}

\begin{command}

>>> csv_path = vna_tools.measure_and_save(znb)

\end{command}

\subsection{現在のNarrow S11から、最小振幅の周波数を取得する}

\begin{command}

>>> f0, amp_min = vna_tools.find_min_freq(znb, 2, "Trc4")

\end{command}

\begin{itemize}

\item TM110：2, "Trc4"

\item TM210：4, "Trc9"

\end{itemize}

\subsection{操作終了後、通信を閉じる}

\begin{command}

>>> znb.close()

>>> rm.close()

\end{command}

\section{Piezo調整}

\subsection{スイッチをVNA側・使用するモードに切り替える}

\begin{command}

% python Switchcontrol.py VNA TM110

\end{command}

\begin{itemize}

\item TM210の場合はTM210を指定する。

\end{itemize}

\subsection{ANC350を初期化する：piezo\_move.py}

\begin{itemize}

\item ANC350をUSBでWindows PCに接続する。

\item Windows PCで piezo\_move.py を実行し、初期化する。

\item USBをWindows PCからLinux PCへ差し替える。

\end{itemize}

\subsection{Linux PCで受信プログラムを起動する：frequency\_controller.py}

\begin{command}

% python frequency_controller.py

\end{command}

\begin{itemize}

\item Waiting for frequency command...と表示されたら、別タブから指示を送る。

\item このターミナルは起動したままにする。

\item この受信側のターミナルに、現在の周波数・目標との差・移動step数・最終結果が表示される。

\item 調整終了後は、次の指示を待ち受ける。

\end{itemize}

\subsection{別タブでPythonを起動し、送信側を準備する}

\subsection{python起動後}

\begin{command}

>>> from frequency_command import FrequencyCommandProducer

>>> producer = FrequencyCommandProducer()

>>> producer.send_target("TM110", 1897400)

\end{command}

\begin{itemize}

\item 引数1：TM110 or TM210

\item 引数2：目標周波数 ［kHz］

\end{itemize}

\subsection{送信側の通信を閉じる}

\begin{command}

>>> producer.close()

\end{command}

\section{Piezoの個別操作}

\subsection{スイッチをVNA側・使用するモードに切り替える}

\begin{command}

% python Switchcontrol.py VNA TM110

\end{command}

\begin{itemize}

\item TM210の場合はTM210を指定する。

\end{itemize}

\subsection{ANC350を初期化する：piezo\_move.py}

\begin{itemize}

\item ANC350をUSBでWindows PCに接続する。

\item Windows PCで piezo\_move.py を実行し、初期化する。

\item USBをWindows PCからLinux PCへ差し替える。

\end{itemize}

\subsection{Linux PCでPythonを起動し、ANC350に接続する}

\begin{command}

>>> from pylablib.devices import Attocube

>>> atc = Attocube.ANC350()

\end{command}

\begin{itemize}

\item frequency\_controller.pyとの同時接続はできない

\end{itemize}

TM110：axis1、TM210：axis2


\subsection{現在位置を確認する}

\begin{command}

>>> atc.get_position(axis)

\end{command}

\begin{itemize}

\item 指定した軸の位置を表示する。単位はm。

\item >>> atc.get\_position() で全軸の確認

\end{itemize}

\subsection{指定した位置へ移動する：絶対位置指定}

\begin{command}

>>> atc.move_to(axis, 0.001)

\end{command}

\begin{itemize}

\item 例：位置 0.001 m へ移動する。

\item 移動量ではなく、移動先の座標を指定する。

\end{itemize}

\subsection{現在位置から指定した距離だけ移動する：相対移動}

\begin{command}

>>> atc.move_by(axis, 10e-6)

\end{command}

\begin{itemize}

\item 例：現在位置から ＋10 µm 移動する。

\item 負方向なら-10e-6を指定する。

\end{itemize}

\subsection{指定したstep数だけ移動する}

\begin{command}

>>> atc.move_by_steps(axis, 10, 0.01)

\end{command}

\begin{itemize}

\item 引数1：軸番号

\item 引数2：移動step数。負方向なら-10などを指定する。

\item 引数3：step間の待ち時間［秒］

\end{itemize}

\subsection{操作終了後、通信を閉じる}

\begin{command}

>>> atc.close()

\end{command}

\section{パワーのfeedback}

\subsection{パワー測定値の送信を開始する：RaS\_NRP2\_Operator.py}

\begin{command}

% python RaS_NRP2_Operator.py

\end{command}

\begin{itemize}

\item このターミナルは起動したままにする。

\end{itemize}

\subsection{スイッチをAMP側・使用するモードに切り替える}

\begin{command}

% python Switchcontrol.py AMP TM110

\end{command}

\begin{itemize}

\item TM210の場合はTM210を指定する。

\end{itemize}

\subsection{別タブでPythonを起動し、フィードバックを準備する}

\subsection{python起動後}

\begin{command}

>>> from power_feedback import PowerFeedbackController

>>> feedback = PowerFeedbackController()

>>> feedback.start(frequency_ghz=1.8974, target_power_w=7e-6, initial_amp_power_dbm=-25.0)

\end{command}

\begin{itemize}

\item frequency\_ghz：AMPの周波数［GHz］

\item target\_power\_w：NRP2 CH2で測定する目標パワー［W］

\item initial\_amp\_power\_dbm：開始時のAMP出力［dBm］

\item 周波数設定 → 初期AMP出力設定 → AMP出力ON → フィードバック開始まで行う。

\end{itemize}

\subsection{設定の変更}

変更後の設定は、次の制御計算から使われる


目標パワーを変更する


\begin{command}

>>> feedback.set_target_power(7e-6)

\end{command}

\begin{itemize}

\item 引数：目標パワー［W］

\end{itemize}

PIDの係数を変更する


\begin{command}

>>> feedback.set_pid(kp=1.0, ki=0.0, kd=0.0)

\end{command}

制御の更新間隔を変更する


\begin{command}

>>> feedback.set_interval(1.0)

\end{command}

\begin{itemize}

\item 引数：更新ごとの待ち時間［秒］

\end{itemize}

\subsection{現在の状態を確認する}

\begin{command}

>>> feedback.get_status()

\end{command}

\begin{itemize}

\item AMP出力ON/OFF、フィードバック状態、周波数、AMP出力、目標値、最新の測定値などを表示する。

\end{itemize}

\subsection{フィードバックを停止し、SG出力をOFFにする}

\begin{command}

>>> feedback.stop()

\end{command}

\subsection{ログのパスを確認}

\begin{command}

>>> print(feedback.log_path)

\end{command}

\subsection{測定ログのリアルタイム確認}

別のタブで


\begin{command}

% tail -f ログのパス

\end{command}

\section{1port運用}

TM110・TM210の両方を、TM110側の入力ポートで測定する。


\subsection{スイッチ操作はSwitchcontrol\_1port.pyを使用する}

\begin{itemize}

\item TM110／TM210のどちらを指定しても、物理入力はTM110側に固定される。

\end{itemize}

\subsection{Piezo調整の受信側はfrequency\_controller\_1port.pyを使用する}

\begin{itemize}

\item 周波数調整の処理・piezo軸・周波数範囲は従来版と同じ

\item 保存するCSV・PNGの名前に1port\_input\_TM110を追加。

\item 保存する図に、入力がTM110側であることを表示。

\end{itemize}

\end{document}
